%=====================================================================
%  CE-1115 Seguridad de la Informacion -- Proyecto, Parte 1
%  Consultas al cliente: preguntas escritas para la reunion con
%  Verdeazul Carbono S.A. (el profesor, en el papel del cliente)
%  Documento de trabajo del Grupo 6. No es uno de los cinco entregables:
%  lo que se responda pasa a la minuta (Anexo A de P1_Politicas).
%
%  Compilar dos veces:  pdflatex P1_ConsultasCliente_grupo6.tex
%=====================================================================
\documentclass[11pt,letterpaper]{article}

\usepackage[utf8]{inputenc}
\usepackage[T1]{fontenc}
\usepackage[spanish,es-noshorthands,es-tabla]{babel}
\usepackage{XCharter}
\usepackage[defaultsans,scale=0.92]{lato}
\usepackage[scaled=0.9]{inconsolata}
\usepackage[letterpaper,margin=2.2cm,headheight=15pt]{geometry}
\usepackage{microtype}
\usepackage{xcolor}
\usepackage{colortbl}
\usepackage{booktabs,tabularx,array,longtable}
\usepackage{enumitem}
\usepackage{titlesec}
\usepackage{fancyhdr}
\usepackage[most]{tcolorbox}
\usepackage{needspace}
\usepackage[hidelinks]{hyperref}

\newcommand{\numgrupo}{6}
\newcommand{\titulocorto}{Consultas al cliente}

\definecolor{vz}{HTML}{0F6E6A}
\definecolor{vzclaro}{HTML}{E7F2F1}
\definecolor{tinta}{HTML}{1A1A1A}
\definecolor{gris}{HTML}{5F5E5A}
\definecolor{regla}{HTML}{B9B8B0}
\definecolor{pend}{HTML}{9A6200}
\definecolor{pendfondo}{HTML}{FFF4D9}
\definecolor{negado}{HTML}{F2F1EC}
\color{tinta}

\titleformat{\section}{\Large\sffamily\bfseries\color{vz}}{\thesection}{0.6em}{}
\titleformat{\subsection}{\large\sffamily\bfseries\color{tinta}}{}{0em}{}
\titlespacing*{\section}{0pt}{1.2em}{0.4em}
\titlespacing*{\subsection}{0pt}{0.9em}{0.2em}
\setlength{\parindent}{0pt}
\setlength{\parskip}{0.5em}
\setlist{leftmargin=1.5em,itemsep=2pt,topsep=3pt,parsep=0pt}

\pagestyle{fancy}
\fancyhf{}
\fancyhead[L]{\small\sffamily\color{gris}CE-1115 $\cdot$ Proyecto Parte 1 $\cdot$ Grupo \numgrupo}
\fancyhead[R]{\small\sffamily\color{gris}\titulocorto}
\fancyfoot[C]{\small\sffamily\color{gris}\thepage}
\renewcommand{\headrulewidth}{0.4pt}
\renewcommand{\headrule}{\hbox to\headwidth{\color{regla}\leaders\hrule height \headrulewidth\hfill}}

\newcolumntype{L}[1]{>{\raggedright\arraybackslash}p{#1}}
\newcolumntype{Y}{>{\raggedright\arraybackslash}X}
\newcommand{\cab}[1]{\textbf{\sffamily #1}}
\newcommand{\porllenar}{\textcolor{pend}{\sffamily\small[por completar]}}
\newcommand{\pri}{\textcolor{vz}{\sffamily\bfseries A}}      % prioridad A
\newcommand{\bloque}[1]{\multicolumn{5}{@{}l}{\cellcolor{vzclaro}\cab{#1}} \\}
\newcommand{\carta}[1]{«\emph{#1}»}

\begin{document}

%--------------------------- encabezado ------------------------------
{\sffamily\large\color{gris} Proyecto $\cdot$ Parte 1 $\cdot$ Diseño del control de acceso}\\[0.25cm]
{\sffamily\bfseries\fontsize{22}{26}\selectfont Consultas al cliente}\\[0.15cm]
{\sffamily\large Preguntas escritas para la reunión con Verdeazul Carbono S.A.}\\[0.2cm]
{\color{vz}\rule{\textwidth}{1.4pt}}

\vspace{0.2cm}
\begin{tabularx}{\textwidth}{@{}L{3.6cm} Y@{}}
\textbf{Equipo consultor} & Grupo \numgrupo: Randall Bolaños López (2019043784), Javier Tenorio Cervantes (2020065308), Julio Varela Venegas (2019008041), Axel Flores Lara (2021453573) \\
\textbf{Parte interesada} & Prof. MSc. Andrés Vargas Rivera, en el papel de Verdeazul Carbono S.A. \\
\textbf{Reunión} & Fecha y hora: \porllenar \quad Medio: \porllenar \\
\end{tabularx}

%=====================================================================
\section{Para qué son estas preguntas}
%=====================================================================

La carta de encargo dice a qué se dedica Verdeazul y qué le preocupa, pero deja abiertos varios puntos
de los que depende el diseño. El equipo ya entregó un diseño completo, construido con supuestos de
trabajo. Esta reunión es para sustituir esos supuestos por lo que el cliente decida.

\textbf{Cómo se usa en la reunión.}

\begin{itemize}
\item Las preguntas están redactadas en el lenguaje del negocio. La columna \emph{Qué decide} es para el
      equipo: dice qué parte del diseño depende de la respuesta.
\item La columna \emph{Si no hay respuesta} es el supuesto con que hoy está hecho el diseño. Si el
      cliente lo confirma, no cambia nada.
\item La prioridad \pri\ marca las doce preguntas que cambian el diseño. Si el tiempo no alcanza, se
      hacen esas primero, en el orden de la sección~\ref{sec:prioridad}.
\item Una persona pregunta y otra anota. Lo respondido se pasa el mismo día a la minuta, que es el
      Anexo~A del informe de políticas.
\end{itemize}

%=====================================================================
\section{Las preguntas}
%=====================================================================

{\small
\renewcommand{\arraystretch}{1.28}
\setlength{\tabcolsep}{4.5pt}
\begin{longtable}{@{}L{0.95cm} L{0.5cm} L{5.5cm} L{4.4cm} L{4.3cm}@{}}
\toprule
\cab{N.º} & \cab{Pr.} & \cab{Pregunta al cliente} & \cab{Qué decide} & \cab{Si no hay respuesta} \\
\midrule
\endfirsthead
\toprule
\cab{N.º} & \cab{Pr.} & \cab{Pregunta al cliente} & \cab{Qué decide} & \cab{Si no hay respuesta} \\
\midrule
\endhead
\bottomrule
\endfoot

\bloque{A. Quién mide y quién aprueba}
P-01 & \pri & ¿Un verificador sale alguna vez a medir, por ejemplo para cubrir a un técnico? ¿Hay verificadores que antes fueron técnicos? &
Si los puestos de técnico y verificador se excluyen siempre o solo dentro de una misma sesión. &
Nunca se mezclan, pero puede haber ascensos. Por eso la regla se comprueba medición por medición. \\
P-02 & \pri & ¿Cuántos técnicos y cuántos verificadores hay? ¿Cuántas fincas son en total y cuántas atiende un técnico en un período? &
Si la separación y la rotación son viables, y los umbrales de consulta. &
Hay al menos dos verificadores. \\
P-03 & B & Cuando una medición trae un error evidente, ¿el verificador lo corrige o siempre la devuelve al técnico? &
Si el verificador puede modificar una medición. &
Siempre la devuelve; nunca la edita. \\
P-04 & B & ¿El certificado se emite solo cuando se aprueba la medición, o alguien más lo firma? ¿Usan hoy firma digital certificada? &
Si la aprobación lleva firma digital personal o un sello de la empresa. &
Lo emite el sistema al aprobar, con un sello de la empresa. \\

\bloque{B. Períodos y enmiendas}
P-05 & \pri & ¿Quién abre y quién cierra un período de medición? ¿Quién decide qué técnico va a cada finca? &
Qué puesto tiene esos permisos. &
La Gerencia de Certificación hace las tres cosas. \\
P-06 & \pri & ¿La revisión de los verificadores termina antes del cierre, o sigue después? ¿Qué pasa con una medición que al cierre no se ha aprobado? &
Los estados del período y qué se puede hacer después de cerrar. &
Lo que no esté aprobado al cierre queda fuera y solo se retoma con una enmienda. \\
P-07 & \pri & Después del cierre, ¿quién puede pedir una corrección y quién debe autorizarla? ¿Debe revisarla un verificador distinto del que aprobó la medición original? &
El flujo de la enmienda y cuántas personas intervienen. &
La pide el técnico o el verificador, la autoriza la Gerencia y la revisa un verificador que no participó. \\
P-08 & \pri & Si la corrección cambia las toneladas de un certificado ya emitido, o ya vendido, ¿qué se hace con ese certificado? &
Si un certificado se puede revocar y reemplazar. &
Se revoca y se emite uno nuevo que hace referencia al anterior. Nada se borra. \\
P-09 & B & ¿Cuántas correcciones de este tipo esperan por período: una o dos, o decenas? &
Cuánta fricción es aceptable en la enmienda. &
Son pocas; se acepta un trámite de tres personas. \\

\bloque{C. El técnico del último párrafo de la carta}
P-10 & \pri & Ante una consulta masiva como la del año pasado, ¿qué quieren que haga el sistema: avisar, pedir una confirmación, o bloquear? ¿A quién debe avisarle? &
La respuesta de los controles de historial. &
Tres escalones: avisa, pide confirmación con motivo y, si sigue, bloquea y avisa a la Gerencia. \\
P-11 & \pri & ¿Un técnico necesita ver fincas que no tiene asignadas? ¿En qué casos? &
Si el alcance del técnico se limita por asignación. &
No. Solo ve las fincas que tiene asignadas en el período abierto. \\
P-12 & B & ¿Cuál es el horario normal de trabajo? ¿Se mide o se registra de noche o en fines de semana? &
Qué se considera fuera de jornada. &
Lunes a viernes, de 6:00 a 19:00, para la oficina. El registro de campo no se restringe. \\
P-13 & B & ¿Cuántas fincas consulta un verificador en un día normal? ¿Y en la última semana de un período? &
Los umbrales de consulta. &
40 fincas por día es normal; 80 es el tope. \\
P-14 & C & ¿El personal firma un acuerdo de confidencialidad? ¿Sabe que sus consultas van a quedar registradas? &
El respaldo legal de los controles. &
Lo firma al ingresar, junto con la política de uso aceptable. \\

\bloque{D. Dueños de finca y auditores}
P-15 & \pri & ¿Los dueños van a entrar ellos mismos al sistema? ¿Un dueño puede tener varias fincas, delegar en un familiar o ser una sociedad? ¿Algún empleado o familiar de empleado es dueño de una finca certificada? &
Cómo se identifica a un dueño y si hay conflictos de interés. &
Entran ellos mismos, pueden autorizar representantes, y una finca de un empleado queda fuera de sus asignaciones. \\
P-16 & B & ¿El dueño debe ver la medición apenas se registra, o solo cuando ya está aprobada? &
Qué mediciones ve el dueño. &
Solo las aprobadas. \\
P-17 & C & ¿Qué datos personales de los dueños guardan hoy? ¿Firman algún consentimiento? &
Las obligaciones de la Ley 8968. &
Identificación y contacto, con consentimiento firmado al darse de alta. \\
P-18 & \pri & ¿Quiénes son los auditores y bajo qué norma auditan? ¿Cuántos días dura la visita, trabajan en sitio o a distancia, y necesitan llevarse copias? &
La cuenta de auditor y qué exige la norma sobre registros. &
Personas con nombre, con acceso de solo lectura durante la visita, sin copias de datos personales. \\
P-19 & B & La trazabilidad de diez años, ¿incluye quién consultó cada dato, o solo quién lo cambió y quién lo aprobó? &
Cuánto tiempo se guardan las consultas. &
Cambios y aprobaciones, diez años. Consultas, dos años. \\

\bloque{E. Administración y operación}
P-20 & B & ¿Quién le pide las altas de personal al administrador y quién las aprueba? ¿Quién lo sustituye cuando no está? &
Quién da la segunda firma en cuentas y roles. &
La Gerencia aprueba todo rol; nadie se asigna permisos a sí mismo. \\
P-21 & B & ¿Qué identifica una hectárea en sus registros: el plano, un polígono, un código? ¿Cómo ocurrió la doble emisión? &
La regla que impide emitir dos veces. &
La geometría de la parcela; el sistema rechaza parcelas traslapadas. \\
P-22 & \pri & ¿Los técnicos registran en la finca, o al volver a la oficina? ¿Hay señal en las fincas? ¿Los equipos son de la empresa? &
Si hace falta registrar sin conexión. &
Registran en línea, con un equipo de la empresa. \\
P-23 & \pri & ¿Dónde estaría el sistema y quién lo operaría? ¿«Solo software libre» deja presupuesto para un servidor y para llaves físicas de seguridad? &
El hospedaje, quién custodia las llaves y el tipo de segundo factor. &
Un servidor operado por un proveedor; sin compra de llaves físicas. \\
P-24 & B & ¿Cuánto tiempo puede estar caído el sistema con un período abierto sin que sea un problema serio? &
Si hace falta un segundo servidor. &
Se tolera un día. No se tolera perder datos. \\
P-25 & C & ¿Alguien de fuera, como un comprador, necesita comprobar que un certificado es auténtico? &
Si hay una consulta pública de certificados. &
Conviene, pero queda fuera de esta fase. \\
P-26 & C & ¿Cuántos años de certificados hay en las hojas de cálculo actuales, y qué debe pasar con ellos? &
Cómo se trata el histórico. &
Se migra una sola vez y queda marcado como histórico. \\
\end{longtable}
}

{\small \textbf{Pr.}: \pri\ cambia el diseño; B ajusta un valor o un detalle; C confirma un supuesto.}

%=====================================================================
\section{Si el tiempo es corto: las doce primero}
\label{sec:prioridad}
%=====================================================================

En este orden, que va de lo que más cambia el diseño a lo que menos:

\begin{enumerate}[itemsep=1pt]
\item \textbf{P-01} y \textbf{P-07}: quién puede medir, aprobar y autorizar una enmienda. De aquí sale la
      técnica de separación de responsabilidades.
\item \textbf{P-10} y \textbf{P-11}: qué hacer con el técnico del último párrafo.
\item \textbf{P-05} y \textbf{P-06}: quién abre y cierra períodos, y qué pasa al cierre.
\item \textbf{P-08}: qué ocurre con un certificado ya emitido cuando hay una enmienda.
\item \textbf{P-02}: cuánta gente hay. Sin ese número no se sabe si la separación es viable.
\item \textbf{P-18} y \textbf{P-15}: auditores y dueños.
\item \textbf{P-22} y \textbf{P-23}: conexión en el campo, hospedaje y presupuesto.
\end{enumerate}

Son las tres que el enunciado señala como abiertas a propósito (P-01, P-07 y P-10) y las nueve que
el equipo encontró al diseñar.

%=====================================================================
\section{Después de la reunión: qué se ajusta}
%=====================================================================

Para el equipo. Cada respuesta que difiera del supuesto obliga a tocar estos lugares. En los cinco
entregables, los puntos pendientes están marcados con la etiqueta de la pregunta y se encuentran
buscando la palabra \texttt{PENDIENTE} o el número, por ejemplo \texttt{P-07}, en el archivo
\texttt{.tex}.

{\small
\renewcommand{\arraystretch}{1.22}
\begin{longtable}{@{}L{1.3cm} L{3.3cm} L{3.8cm} L{3.4cm} L{3.4cm}@{}}
\toprule
\cab{Pregunta} & \cab{Políticas (E1)} & \cab{Control de acceso (E2)} & \cab{Matriz (E3)} & \cab{Arquitectura (E4)} \\
\midrule
\endfirsthead
\toprule
\cab{Pregunta} & \cab{Políticas (E1)} & \cab{Control de acceso (E2)} & \cab{Matriz (E3)} & \cab{Arquitectura (E4)} \\
\midrule
\endhead
\bottomrule
\endfoot
P-01 & --- & §3.4, §5.2, §5.4 & Exclusión TEC--VER (§4) & --- \\
P-02 & --- & §3.4, §5.3 & Cardinalidad (§4) & --- \\
P-03 & --- & §5.5 & (VER, Mediciones) & --- \\
P-04 & §5, Ley 8454 & --- & --- & Anexo B; emisión \\
P-05 & --- & §3.1 & Columnas Períodos y Asignaciones & --- \\
P-06 & --- & --- & C3, C14 & Estados (Anexo A.2) \\
P-07 & --- & §4.1 (H4), §5.5 & C8, C10 & --- \\
P-08 & --- & --- & C4 & Anexo B; estados \\
P-09 & --- & §5.5 & --- & --- \\
P-10 & --- & §4.1 (H1, H3), §4.2 & H1, H3 & --- \\
P-11 & --- & §4 & C1 & --- \\
P-12, P-13 & --- & §4.1 (valores iniciales) & H1, H2 & --- \\
P-14 & §5, POL-CA-02 & §4.2 & --- & --- \\
P-15 & POL-CA-08 & --- & §4, C5 & --- \\
P-16 & --- & --- & C5 & --- \\
P-17 & §2.2 (A-04), §5 & --- & --- & --- \\
P-18 & §5, POL-CA-07 & --- & C6 & --- \\
P-19 & POL-BT-04 & §4.1 & --- & §4.5 \\
P-20 & --- & §3.4 & C12 & --- \\
P-21 & --- & --- & C11 & §4.2 \\
P-22 & POL-AU-11 & --- & --- & RR-8; Anexo B \\
P-23 & POL-CA-09, POL-AU-05 & --- & --- & §1, RR-7 \\
P-24 & §3 & --- & --- & RR-6 \\
P-25 & AM-10 & --- & --- & Anexo B \\
P-26 & §6.4 & --- & --- & RR-9 \\
\end{longtable}
}

Además, en todos los casos: llenar el Anexo~A y actualizar el Anexo~B del informe de políticas, quitar
los bloques de pendiente de los cinco documentos, completar los dos bloques del documento del atributo
DI, actualizar la lámina de pendientes de la presentación y recompilar.

%=====================================================================
\section{Consultas de logística al profesor}
%=====================================================================

Fuera del papel de cliente, conviene confirmar tres cosas del curso:

\begin{enumerate}[itemsep=1pt]
\item ¿Cuánto dura la defensa oral y cuánto de ese tiempo es para preguntas?
\item ¿La minuta tiene un formato obligatorio, o basta la plantilla del Anexo~A?
\item ¿El número de grupo que va en el asunto del correo y en los nombres de archivo es el \numgrupo?
\end{enumerate}

\end{document}
